% docs/manual/latex/user/09_performance.tex - chapter of the LaTeX edition of the book "user_manual" (SIMC++ documentation).
%   Converted once from docs/manual/markdown/user/09_performance.md by tools/md2latex.py; from then on THIS FILE is the source of the
%   LaTeX edition and may be edited by hand. Re-running tools/md2latex.py --write-sources for this
%   chapter OVERWRITES the hand edits (see docs/README.md).
%   Compile: cmake --build <build> --target pdf, or tools/build_docs_pdf.py (figures, xelatex passes,
%   makeindex with docs/common/latex/simcpp.ist); by hand, see docs/README.md. XeLaTeX, class
%   docs/common/latex/simcpp.cls (packages: fontspec, tikz, xcolor, hyperref, listings, longtable, booktabs,
%   titlesec, titletoc, fancyhdr, makeidx, preview, ...); fonts TeX Gyre Pagella, TeX Gyre Heros, DejaVu Sans
%   Mono. The figures are built first from docs/{manual,translation}/figures/*.tex (preview.sty) into PDFs.
\chapter{\texorpdfstring{Performance tips}{Performance tips}}\label{user_manual-09u5fperformance}
\simcabstract{Measured on the development machine (Intel i7-3770, 4 cores / 8 hardware threads, GCC 16.2 \code{-{}O3}; \code{docs/\allowbreak{}manual/\allowbreak{}markdown/\allowbreak{}reference/\allowbreak{}PERFORMANCE.\allowbreak{}md}).}
\section{\texorpdfstring{Speed}{Speed}}\label{user_manual-09u5fperformance:91-speed}
One thread, whole program including the start-up of about 0.5 s:

\begin{simctable}{small}{2}{>{\raggedright\arraybackslash}p{\dimexpr 0.272\linewidth-2\tabcolsep\relax}>{\raggedright\arraybackslash}p{\dimexpr 0.136\linewidth-2\tabcolsep\relax}}
%
\toprule\rowcolor{simchead}\simcth{case} & \simcth{µs per try} \\\midrule\endfirsthead
\multicolumn{2}{@{}l}{\simccontinued}\\\toprule\rowcolor{simchead}\simcth{case} & \simcth{µs per try} \\\midrule\endhead
\midrule\multicolumn{2}{r@{}}{\simccontinues}\\\endfoot
\bottomrule\endlastfoot
H(e,e'p) radiative & 16.3 \\
π+ H, decay on & 27.1 \\
SIDIS & 7.2 \\
rho & 14.2 \\
π0 calorimeter (NPS) & 5.8 \\
polarized-target field & 40.7 \\
\end{simctable}
The cost of a run is (number of tries) × (cost per
try). The number of tries per accepted event is Ntried/\allowbreak{}Ncontribute in the \code{.\allowbreak{}hist}: about 21 for the H(e,e'p) example,
about 3000 for rho.

\section{\texorpdfstring{Tips}{Tips}}\label{user_manual-09u5fperformance:92-tips}
\begin{itemize}
\item \textbf{Use threads.} \code{-{}-{}threads\ 4}\simcidx{Command-line options!threads@\code{-{}-{}threads}} gave 2.8-3.4× and \code{-{}-{}threads\ 8}\simcidx{Command-line options!threads@\code{-{}-{}threads}} 2.8-4.3× on the 4-core machine (\code{docs/\allowbreak{}manual/\allowbreak{}markdown/\allowbreak{}reference/\allowbreak{}PERFORMANCE.\allowbreak{}md}
§3). Beyond the number of physical cores only hyper-threading helps. Threads are not allowed with \code{-{}-{}fortran-{}compat}\simcidx{Command-line options!fortran-compat@\code{-{}-{}fortran-{}compat}} or
RANLUX (\hyperref[user_manual-08u5frunningu5fandu5freproducibility:85-threads---threads-n]{8.5}).
\item \textbf{Memory with threads.} Each thread holds its own tables; a pion run on 8 threads uses about 450 MB. Each thread also
builds its stages at the start (0.1-0.3 s per thread), which dominates very short runs.
\item \textbf{Set tight \code{SPedge} limits.} The generation volume comes from them, plus the slop. Wide-open defaults (±1e10) waste
almost every try, and the run may never end (\hyperref[user_manual-10u5ffaq:the-run-never-ends]{10}).
\item \textbf{Test with negative \hyperref[param:ngen]{\code{ngen}}\simcidx{Parameters!ngen@\code{ngen}}.} \code{-{}-{}ngen\ -{}100000}\simcidx{Command-line options!ngen@\code{-{}-{}ngen}} makes exactly 100000 tries, whatever the acceptance.
\item \textbf{mt64 is faster than RANLUX} (by about 25 \%). Use RANLUX only to reproduce legacy SIMC results.
\item \textbf{Leave out outputs you do not need.} For example, a \code{pipeline.\allowbreak{}sinks} list without \code{gen\_\allowbreak{}hist} (or \code{Nntu\ =\ 0} for no tree).
The \code{.\allowbreak{}hist} is written by \code{fortran\_\allowbreak{}hist}.
\item \textbf{Do not build with \code{-{}march=\allowbreak{}native} or \code{-{}ffast-{}math}.} They change the floating-point results; the bitwise tests and
the exact reproduction of legacy results then fail.
\item Radiative runs spend most of their time in the radiative-correction set-up of each try (the gamma function), the
target energy loss and the COSY maps. Runs without radiation are faster.
\end{itemize}
